\documentclass[10pt,letterpaper]{article}

\usepackage[top=0.4in,bottom=0.4in,left=0.5in,right=0.5in]{geometry}
\usepackage[utf8]{inputenc}
\usepackage[T1]{fontenc}
\usepackage[spanish]{babel}
\usepackage{enumitem}
\usepackage[hidelinks]{hyperref}
\usepackage{titlesec}
\usepackage{xcolor}

\definecolor{heading}{HTML}{1F2937}

\pagestyle{empty}
\setlength{\parindent}{0pt}
\setlength{\parskip}{0pt}
\setlist[itemize]{leftmargin=*, itemsep=0pt, topsep=1pt, parsep=0pt, partopsep=0pt}

\titleformat{\section}
  {\large\bfseries\color{heading}}
  {}{0pt}{}
  [\vspace{-3pt}\titlerule]
\titlespacing*{\section}{0pt}{3pt}{1pt}

% Proyecto: nombre -- fecha (derecha) -- stack en segunda l\'inea
\newcommand{\project}[3]{\vspace{2pt}\textbf{#1} \hfill #2 \\ \textit{#3} \par\vspace{1pt}}

% Experiencia: org, fechas, rol
\newcommand{\resumeSubheading}[3]{%
  \vspace{2pt}\textbf{#1} \hfill #2 \\
  \textit{#3} \\
}

\begin{document}

\begin{center}
    {\LARGE \textbf{Carolina de Jes\'us Ortega Cepeda}} \\
    \vspace{2pt}
    Ingenier\'ia en Tecnolog\'ias Computacionales \\
    \vspace{3pt}
    \href{tel:8124300714}{+52 812 430 0714} $|$
    \href{mailto:carodejesus100@gmail.com}{carodejesus100@gmail.com} $|$
    \href{https://www.linkedin.com/in/caro-ortega-cepeda-83ba3424b}{LinkedIn} $|$
    % TODO: Confirmar usuario/URL de GitHub.
    \href{https://github.com/AthensCort}{GitHub} $|$
    Monterrey, M\'exico
\end{center}

\section{Perfil}
Desarrolladora de Software (Intern) en ciberseguridad y estudiante de ITC que construye productos completos en Python y TypeScript. Dise\~no APIs REST, modelos de datos relacionales (PostgreSQL) y frontends en React, con enfoque en herramientas de seguridad. Despliego apps en producci\'on de forma independiente y en equipos con flujos basados en Git.

\section{Habilidades T\'ecnicas}
\textbf{Lenguajes:} Python, TypeScript, JavaScript, SQL, Java \\
\textbf{Frontend:} React, TailwindCSS, HTML/CSS \\
\textbf{Backend:} Node.js, Express, Flask, APIs REST, autenticaci\'on JWT \\
\textbf{Bases de datos:} PostgreSQL, SQL Server, Prisma ORM \\
\textbf{DevOps / Herramientas:} Docker, Git/GitHub, Vercel

\section{Experiencia}
\resumeSubheading{Mosaic Holdings}{Jun 2025 -- Actualidad}{Desarrolladora de Software (Intern) -- Ciberseguridad}
\begin{itemize}
    \item Desarrollo software interno de seguridad para el equipo de ciberseguridad de la empresa, reportando directamente al gerente de ciberseguridad.
    \item Dise\~n\'e y constru\'i \textbf{Raccoon Scout}, una plataforma full-stack de detecci\'on de phishing/typosquatting (Python/Flask, React/TypeScript) que genera dominios similares con \texttt{dnstwist} y resuelve datos de DNS (A/NS/MX), WHOIS y GeoIP.
    \item Implement\'e una API REST en Flask con trabajos de escaneo as\'incronos y multihilo (ThreadPoolExecutor) y endpoints de sondeo; persist\'i resultados en Supabase/PostgreSQL con RLS, Google OAuth y respaldo en SQLite.
    \item Constru\'i el frontend en React 19 con mapa mundial 3D interactivo (Three.js, TopoJSON); despliegue en Railway (Gunicorn) + Vercel, contenerizado con Docker.
\end{itemize}

\resumeSubheading{VOLTEC -- Robótica FIRST Tech Challenge}{01/2024 -- 06/2024}{Mentora T\'ecnica}
\begin{itemize}
    \item Depur\'e y mejor\'e el c\'odigo de control de robots, gestion\'e la log\'istica de desarrollo y operaci\'on, y mantuve documentaci\'on t\'ecnica sobre programaci\'on y hardware.
    \item Apoy\'e la clasificaci\'on nacional del equipo en el Campeonato FTC 2024.
\end{itemize}

\section{Proyectos}
\project{MyoQuack -- Sistema de Biofeedback y Rehabilitaci\'on EMG}{Mar -- Abr 2026}{TypeScript, Node.js, Express, Prisma, PostgreSQL, ESP32 \quad $\cdot$\ Individual}
\begin{itemize}
    \item Constru\'i un flujo completo de IoT a web: firmware en ESP32 que captura se\~nales musculares (EMG), un backend REST tipado en Node.js/Express + Prisma sobre PostgreSQL, y un cliente en React.
    \item Model\'e m\'edicos, pacientes, sesiones EMG y eventos por contracci\'on (pico $\mu$V, RMS, iEMG, waveform); implement\'e autenticaci\'on JWT + bcrypt, acceso por rol, exportaci\'on a CSV, documentaci\'on Swagger/OpenAPI y Docker Compose.
\end{itemize}

\project{Ternium -- Sistema iOS de Gesti\'on de Turnos y Ventanillas}{Ago -- Dic 2025}{Swift, SwiftUI, MVVM, Flask, SQL Server \quad $\cdot$\ Equipo}
\begin{itemize}
    \item Constru\'i dos apps nativas de iOS (SwiftUI, MVVM): una app cliente para generar, seguir y cancelar turnos, y una app de administraci\'on con autenticaci\'on por rol, estad\'isticas de empleados y dashboards de ventanillas.
    \item Dise\~n\'e un backend en Flask + SQL Server que expone la API REST de turnos/auth/dashboard, consumida desde iOS v\'ia URLSession y DTOs tipados, con manejo de TLS/transporte para actualizaciones en vivo.
\end{itemize}

\project{AFI -- Plataforma de Interacci\'on con Aficionados}{Feb -- Jun 2026}{React 19, TypeScript, Vite, TailwindCSS, Supabase \quad $\cdot$\ Equipo}
\begin{itemize}
    \item Desarroll\'e el m\'odulo de salas de chat en tiempo real (crear/unirse a salas, mensajer\'ia en vivo) con React, hooks personalizados y Supabase Realtime, m\'as el sistema de amigos (invitaciones, notificaciones) y funciones de perfil/avatar.
    \item Escrib\'i migraciones SQL para solicitudes de uni\'on a salas y relaciones de amistad; contribu\'i a la tienda por puntos con registro de auditor\'ia para administradores.
\end{itemize}

% Eliminado para caber en una p\'agina -- reactivar si se necesita:
% \project{Simulaci\'on Multi-Agente en Cuadr\'icula}{Ago -- Dic 2025}{Python, AgentPy, NumPy, Matplotlib, Seaborn}
% \begin{itemize}
%     % TODO: Confirmar si debe describirse como un modelo de cruce de tr\'afico (el c\'odigo es un agente de limpieza en cuadr\'icula).
%     \item Constru\'i un modelo basado en agentes de un agente aut\'onomo que barre una cuadr\'icula (cobertura bustrof\'edica), registrando estado, conteo de movimientos y \% de cobertura, con visualizaciones animadas.
% \end{itemize}

\section{Educaci\'on}
\resumeSubheading{Tecnol\'ogico de Monterrey, Campus Monterrey}{Graduaci\'on estimada: junio 2027}{Ingenier\'ia en Tecnolog\'ias Computacionales \quad Promedio: 96.38/100}
\textbf{Materias relevantes:} Estructuras de Datos y Algoritmos, An\'alisis y Dise\~no de Algoritmos Avanzados, An\'alisis de Requerimientos de Software, Programaci\'on Orientada a Objetos, Seguridad Computacional en Redes y Sistemas de Software, Internet de las Cosas.

\end{document}
